\section{AIME I 1983}
        \begin{enumerate}
        	\item (\textit{AIME I 1983}) Let $x$, $y$ and $z$ all exceed $1$ and let $w$ be a positive number such that $\log_xw=24$, $\log_y w = 40$ and $\log_{xyz}w=12$. Find $\log_zw$.


 
 \textbf{Answer: }060



\textbf{Solution 1}

Converting the given equations to exponential form yields $x^{24}=w$, $y^{40}=w$, and $(xyz)^{12}=w$. 

 Next, converting these new equations so that the left side of each equation is a $120$th power yields $x^{120}=w^5$, $y^{120}=w^3$, and $(xyz)^{120}=x^{120}y^{120}z^{120} = w^{10}$.

 Substituting for $x^{120}$ and $y^{120}$ yields $w^5w^3z^{120}=w^8z^{120}=w^{10}$, so $z^{120}=w^2$.  Therefore, $w = z^{60}$ and $\log_z w=\boxed{060}$.

 





\textbf{Solution 2}

Converting the given equations to exponential form yields. $x^{24}=w$, $y^{40}=w$, and $(xyz)^{12}=x^{12}y^{12}z^{12} = w$. 

 Since $x^{24} = w$, $x = w^{\frac{1}{24}}$ and $x^{12} = w^{\frac{1}{2}}$.  Also, since $y^{40} = w$, $y = w^{\frac{1}{40}}$ and $y^{12} = w^{\frac{3}{10}}$.

 Substituting into $x^{12}y^{12}z^{12} = w$ yields $w^{\frac{1}{2}}w^{\frac{3}{10}}z^{12}=w^{\frac{4}{5}}z^{12}=w$. So $z^{12} = w^\frac{1}{5}$ and $w = z^{60}$. Therefore, $\log_{z} w =\boxed{060}$.

 





\textbf{Solution 3}

Applying the change of base formula,
 \begin{align*} \log_x w = 24 &\implies \frac{\log w}{\log x} = 24 \implies \frac{\log x}{\log w} = \frac 1 {24} \\ \log_y w = 40 &\implies \frac{\log w}{\log y} = 40 \implies \frac{\log y}{\log w} = \frac 1 {40} \\ \log_{xyz} w = 12 &\implies \frac{\log {w}}{\log {xyz}} = 12 \implies \frac{\log x +\log y + \log z}{\log w} = \frac 1 {12} \end{align*}
Therefore, $\frac {\log z}{\log w} = \frac 1 {12} - \frac 1 {24} - \frac 1{40} = \frac 1 {60}$. 

 Hence, $\log_z w = \boxed{060}$.

 





\textbf{Solution 4}

Since $\log_a b = \frac{1}{\log_b a}$, the given conditions can be rewritten as $\log_w x = \frac{1}{24}$, $\log_w y = \frac{1}{40}$, and $\log_w xyz = \frac{1}{12}$. Since $\log_a \frac{b}{c} = \log_a b - \log_a c$,  $\log_w z = \log_w xyz - \log_w x - \log_w y = \frac{1}{12}-\frac{1}{24}-\frac{1}{40}=\frac{1}{60}$. Therefore, $\log_z w = \boxed{060}$.

 





\textbf{Solution 5}

If we convert all of the equations into exponential form, we receive $x^{24}=w$, $y^{40}=w$, and $(xyz)^{12}=w$. The last equation can also be written as $x^{12}y^{12}z^{12}=w$. Also note that by multiplying the first two equations, we get, $x^{24}y^{40}= w^{2}$. Taking the square root of this, we find that $x^{12}y^{20}=w$. Recall, $x^{12}y^{12}z^{12}=w$. Thus, $z^{12}= y^{8}$. Also recall, $y^{40}=w$. Therefore, $z^{60}$ = $y^{40}$ = $w$. So, $\log_z w$ = $\boxed{060}$.

 -Dhillonr25, Bobbob

 





\textbf{Solution 6}

Converting all of the logarithms to exponentials gives $x^{24} = w, y^{40} =w,$ and $x^{12}y^{12}z^{12}=w.$ Thus, we have $y^{40} = x^{24} \Rightarrow z^3=y^2.$ We are looking for $\log_z w,$ which by substitution, is $\log_{y^{\frac{2}{3}}} y^{40} = 40 \div \frac{2}{3} = \boxed{060}.$

 




	\item (\textit{AIME I 1983}) Let $f(x)=|x-p|+|x-15|+|x-p-15|$, where $0 < p < 15$. Determine the \href{https://artofproblemsolving.com/wiki/index.php?title=Minimum}{minimum} value taken by $f(x)$ for $x$ in the \href{https://artofproblemsolving.com/wiki/index.php?title=Interval}{interval} $p \leq x\leq15$.


 
 \textbf{Answer: }015



\textbf{Solution 1}

It is best to get rid of the \href{https://artofproblemsolving.com/wiki/index.php?title=Absolute\_value}{absolute values} first. 

 Under the given circumstances, we notice that $|x-p|=x-p$, $|x-15|=15-x$, and $|x-p-15|=15+p-x$.

 Adding these together, we find that the sum is equal to $30-x$, which attains its minimum value (on the given interval $p \leq x \leq 15$) when $x=15$, giving a minimum of $\boxed{015}$.

 





\textbf{Solution 2}

Let $p$ be equal to $15 - \varepsilon$, where $\varepsilon$ is an almost neglectable value. Because of the small value $\varepsilon$, the \href{https://artofproblemsolving.com/wiki/index.php?title=Domain}{domain} of $f(x)$ is basically the \href{https://artofproblemsolving.com/wiki/index.php?title=Set}{set} ${15}$. plugging in $15$ gives $\varepsilon + 0 + 15 - \varepsilon$, or $15$, so the answer is $\boxed{015}$

 




	\item (\textit{AIME I 1983}) What is the product of the real roots of the equation $x^2 + 18x + 30 = 2 \sqrt{x^2 + 18x + 45}$?


 
 \textbf{Answer: }020



\textbf{Solution 1}

Let $y=x^2+18x+30$. Substituting into the original equation yields $y=2\sqrt{y+15}$.

 Squaring both sides yields $y^2 = 4y + 60$, which can be rearranged to form $y^2 - 4y - 60 = (y-10)(y+6) = 0$.  Though this quadratic has roots $-6$ and $10$, only $y=10$ satisfies the original equation.

 So $y=x^2+18x+30=10$, and therefore $x^2+18x+20=0$.  This quadratic has discriminant $18^2 - 4 \cdot 1 \cdot 20 = 244 > 0$, so both its roots are real. By \href{https://artofproblemsolving.com/wiki/index.php?title=Vieta\%27s\_Formulas}{Vieta's Formulas}, the product of these roots is $\boxed{020}$.

 





\textbf{Solution 2}

We begin by noticing that the polynomial on the left is $15$ less than the polynomial under the radical sign. Thus: 
 \[(x^2+ 18x + 45) - 2\sqrt{x^2+18x+45} - 15 = 0.\] Letting $n = \sqrt{x^2+18x+45}$, we have $n^2-2n-15 = 0 \Longrightarrow (n-5)(n+3) = 0$. Because the square root of a real number can't be negative, the only possible $n$ is $5$.

 Substituting that in, we have 
 \[\sqrt{x^2+18x+45} = 5 \Longrightarrow x^2 + 18x + 45 = 25 \Longrightarrow x^2+18x+20=0.\]

 Reasoning as in Solution 1, the product of the roots is $\boxed{020}$.

 





\textbf{Solution 3}

Begin by completing the square on both sides of the equation, which gives 
 \[(x+9)^2-51=2\sqrt{(x+3)(x+15)}\]Now by substituting $y=x+9$, we get $y^2-51=2\sqrt{(y-6)(y+6)}$, or 
 \[y^4-106y^2+2745=0\]The solutions in $y$ are then 
 \[y=x+9=\pm3\sqrt{5},\pm\sqrt{61}\]Turns out, $\pm3\sqrt{5}$ are a pair of extraneous solutions. Thus, our answer is then 
 \[\left(\sqrt{61}-9\right)\left(-\sqrt{61}-9\right)=81-61=\boxed{020}\]By difference of squares.

 





\textbf{Solution 4}

We are given the equation
 \[x^2+18x+30=2\sqrt{x^2+18x+45}\]Squaring both sides yields
 \[(x^2+18x+30)^2=4(x^2+18x+45)\]
 \[(x^2+18x+30)^2=4(x^2+18x+30+15)\]
 \[(x^2+18x+30)^2=4(x^2+18x+30)+60\]
 \[(x^2+18x+30)^2-4(x^2+18x+30)-60=0\]Substituting $y=x^2+18x+30$ yields
 \[y^2-4y-60=0\]
 \[(y+6)(y-10)=0\]Thus $y=x^2+18x+30=-6,10$. However if $y=-6$, the left side of the equation 
 \[x^2+18x+30=2\sqrt{x^2+18x+45}\]would be negative while the right side is negative. Thus $y=10$ is the only possible value and we have
 \[x^2+18x+30=10\]
 \[x^2+18x+20=0\]Since the discriminant $\sqrt{18^2-4\cdot20}$ is real, both the roots are real. Thus by Vieta's Formulas, the product of the roots is the constant, $\boxed{020}$.

 




	\item (\textit{AIME I 1983}) A machine-shop cutting tool has the shape of a notched circle, as shown. The radius of the circle is $\sqrt{50}$ cm, the length of $AB$ is $6$ cm and that of $BC$ is $2$ cm. The angle $ABC$ is a right angle. Find the square of the distance (in centimeters) from $B$ to the center of the circle.


 
                    \begin{figure}[H]
                        \centering
                        \includegraphics[width=0.5\textwidth]{Problem/AIME_1983_I_4_Problem_Figure_1.png}
                    \end{figure}
                    


 
 \textbf{Answer: }026



\textbf{Solution 1}

Because we are given a right angle, we look for ways to apply the \href{https://artofproblemsolving.com/wiki/index.php?title=Pythagorean\_Theorem}{Pythagorean Theorem}.  Let the foot of the \href{https://artofproblemsolving.com/wiki/index.php?title=Perpendicular}{perpendicular} from $O$ to $AB$ be $D$ and let the foot of the perpendicular from $O$ to the \href{https://artofproblemsolving.com/wiki/index.php?title=Line}{line} $BC$ be $E$. Let $OE=x$ and $OD=y$. We're trying to find $x^2+y^2$.

 
                    \begin{figure}[H]
                        \centering
                        \includegraphics[width=0.5\textwidth]{Solution/AIME_1983_I_4_Solution_Figure_1.png}
                    \end{figure}
                    Applying the Pythagorean Theorem, $OA^2 = OD^2 + AD^2$ and $OC^2 = EC^2 + EO^2$.

 Thus, $\left(\sqrt{50}\right)^2 = y^2 + (6-x)^2$, and $\left(\sqrt{50}\right)^2 = x^2 + (y+2)^2$. We solve this system to get $x = 1$ and $y = 5$, such that the answer is $1^2 + 5^2 = \boxed{026}$.

 





\textbf{Solution 2 (Trig)}

We'll use the \href{https://artofproblemsolving.com/wiki/index.php?title=Law\_of\_cosines}{law of cosines}.  Let $O$ be the center of the circle; we wish to find $OB$.  We know how long $OA$ and $AB$ are, so if we can find $\cos \angle OAB$, we'll be in good shape.

 
                    \begin{figure}[H]
                        \centering
                        \includegraphics[width=0.5\textwidth]{Solution/AIME_1983_I_4_Solution_Figure_1.png}
                    \end{figure}
                    We can find $\cos \angle OAB$ using angles $OAC$ and $BAC$.  First we note that by \href{https://artofproblemsolving.com/wiki/index.php?title=Pythagorean\_theorem}{ Pythagoras},
 \[AC = \sqrt{AB^2 + BC^2} = \sqrt{36 + 4} = \sqrt{40} = 2 \sqrt{10}.\]If we let $M$ be the midpoint of $AC$, that mean that $AM = \sqrt{10}$.  Since $\triangle OAC$ is isosceles ($OA = OC$ from the definition of a circle), $M$ is also the foot of the altitude from $O$ to $AC.$

 
                    \begin{figure}[H]
                        \centering
                        \includegraphics[width=0.5\textwidth]{Solution/AIME_1983_I_4_Solution_Figure_1.png}
                    \end{figure}
                    It follows that $OM = \sqrt{40} = 2 \sqrt{10}$.  Therefore
 \begin{align*} \cos \angle OAC = \frac{\sqrt{10}}{\sqrt{50}} &= \frac{1}{\sqrt{5}}, \\ \sin \angle OAC = \frac{2 \sqrt{10}}{\sqrt{50}} &= \frac{2}{\sqrt{5}}. \end{align*}
Meanwhile, from right triangle $ABC,$ we have
 \begin{align*} \cos \angle BAC = \frac{6}{\sqrt{40}} &= \frac{3}{\sqrt{10}}, \\ \sin \angle BAC = \frac{2}{\sqrt{40}} &= \frac{1}{\sqrt{10}}. \end{align*}


 This means that by the \href{https://artofproblemsolving.com/wiki/index.php?title=Trigonometric\_identities\#Angle\_addition\_identities}{ angle subtraction formulas},
 \begin{align*} \cos \angle OAB &= \cos (\angle OAC - \angle BAC) \\ &= \cos \angle OAC \cos \angle BAC + \sin \angle OAC \sin \angle BAC \\ &= \frac{1}{\sqrt{5}} \cdot \frac{3}{\sqrt{10}} + \frac{2}{\sqrt{5}} \cdot \frac{1}{\sqrt{10}} \\ &= \frac{5}{5 \sqrt{2}} = \frac{1}{\sqrt{2}}. \end{align*}


 Now we have all we need to use the law of cosines on $\triangle OAB.$  This tells us that
 \begin{align*} OB^2 &= AO^2 + AB^2 - 2 AO \cdot AB \cdot \cos \angle OAB \\ &= 50 + 36 - 2 \cdot 5 \sqrt{2} \cdot 6 \cdot \frac{1}{\sqrt{2}} \\ &= 86 - 2 \cdot 5 \cdot 6 \\ &= 26. \end{align*}


 





\textbf{Solution 3}

Mark the midpoint $M$ of $AC$. Then, drop perpendiculars from $O$ to $AB$ (with foot $T_1$), $M$ to $OT_1$ (with foot $T_2$), and $M$ to $AB$ (with foot $T_3$).

 
                    \begin{figure}[H]
                        \centering
                        \includegraphics[width=0.5\textwidth]{Solution/AIME_1983_I_4_Solution_Figure_1.png}
                    \end{figure}
                    First notice that by computation, $OAC$ is a $\sqrt {50} - \sqrt {40} - \sqrt {50}$ isosceles triangle, so $AC = MO$.Then, notice that $\angle MOT_2 = \angle T_3MO = \angle BAC$. Therefore, the two blue triangles are congruent, from which we deduce $MT_2 = 2$ and $OT_2 = 6$. As $T_3B = 3$ and $MT_3 = 1$, we subtract and get $OT_1 = 5,T_1B = 1$. Then the Pythagorean Theorem tells us that $OB^2 = \boxed{026}$.

 





\textbf{Solution 4 (Trig)}

Draw segment $OB$ with length $x$, and draw radius $OQ$ such that $OQ$ bisects chord $AC$ at point $M$. This also means that $OQ$ is perpendicular to $AC$. By the Pythagorean Theorem, we get that $AC=\sqrt{(BC)^2+(AB)^2}=2\sqrt{10}$, and therefore $AM=\sqrt{10}$. Also by the Pythagorean theorem, we can find that $OM=\sqrt{50-10}=2\sqrt{10}$.

 Next, find $\angle BAC=\arctan{\left(\frac{2}{6}\right)}$ and $\angle OAM=\arctan{\left(\frac{2\sqrt{10}}{\sqrt{10}}\right)}$. Since $\angle OAB=\angle OAM-\angle BAC$, we get 
 \[\angle OAB=\arctan{2}-\arctan{\frac{1}{3}}\]
 \[\tan{(\angle OAB)}=\tan{(\arctan{2}-\arctan{\frac{1}{3}})}\]By the subtraction formula for $\tan$, we get
 \[\tan{(\angle OAB)}=\frac{2-\frac{1}{3}}{1+2\cdot \frac{1}{3}}\]
 \[\tan{(\angle OAB)}=1\]
 \[\cos{(\angle OAB)}=\frac{1}{\sqrt{2}}\]Finally, by the Law of Cosines on $\triangle OAB$, we get 
 \[x^2=50+36-2(6)\sqrt{50}\frac{1}{\sqrt{2}}\]
 \[x^2=\boxed{026}.\]

 





\textbf{Solution 5}

We use coordinates. Let the circle have center $(0,0)$ and radius $\sqrt{50}$; this circle has equation $x^2 + y^2 = 50$. Let the coordinates of $B$ be $(a,b)$. We want to find $a^2 + b^2$. $A$ and $C$ with coordinates $(a,b+6)$ and $(a+2,b)$, respectively, both lie on the circle. From this we obtain the system of equations

 $a^2 + (b+6)^2 = 50$

 $(a+2)^2 + b^2 = 50$

 After expanding these terms, we notice by subtracting the first and second equations, we can cancel out $a^2$ and $b^2$. after substituting $a=3b+8$ and plugging back in, we realize that $(a,b)=(-7,-5)$ or $(5,-1)$. Since the first point is out of the circle, we find that $(5,-1)$ is the only relevant answer. This paragraph is written by ~hastapasta. 

 Solving, we get $a=5$ and $b=-1$, so the distance is $a^2 + b^2 = \boxed{026}$.

 





\textbf{Solution 6 (Trigonometry)}


                    \begin{figure}[H]
                        \centering
                        \includegraphics[width=0.5\textwidth]{Solution/AIME_1983_I_4_Solution_Figure_1.png}
                    \end{figure}
                    I will use the \href{https://artofproblemsolving.com/wiki/index.php?title=Law\_of\_cosines}{law of cosines} in triangle $\triangle OAC$ and $\triangle OBC$.

 $AC = \sqrt{AB^2 + BC^2} = \sqrt{6^2 + 2^2} = 2 \sqrt{10}$

 $\cos \angle ACB = \frac{2}{2\sqrt{10}} = \frac{1}{\sqrt{10}}$

 $\cos \angle ACO = \frac{AC^2+OC^2-OA^2}{2 \cdot AC \cdot OC} = \frac{(2\sqrt{10})^2+(\sqrt{50})^2-(\sqrt{50})^2}{2 \cdot 2\sqrt{10} \cdot \sqrt{50}} = \frac{1}{\sqrt{5}}$

 $\sin \angle ACB = \sqrt{1-\cos^2 \angle ACB} = \sqrt{1-(\frac{1}{\sqrt{10}})^2} = \frac{3}{\sqrt{10}}$

 $\sin \angle ACO = \sqrt{1-\cos^2 \angle ACO} = \sqrt{1-(\frac{1}{\sqrt{5}})^2} = \frac{2}{\sqrt{5}}$

 $\cos \angle OCB = \cos (\angle ACB - \angle ACO) = \cos \angle ACB \cdot \cos \angle ACO + \sin \angle ACB \cdot \sin \angle ACO = \frac{1}{\sqrt{10}} \cdot \frac{1}{\sqrt{5}} + \frac{3}{\sqrt{10}} \cdot \frac{2}{\sqrt{5}} = \frac{7}{5\sqrt{2}}$

 $OB^2 = OC^2 + BC^2 - 2 \cdot OC \cdot BC \cdot \cos \angle OCB = (\sqrt{50})^2 + 2^2 - 2 \cdot \sqrt{50} \cdot 2 \cdot \frac{7}{5\sqrt{2}} = 50 + 4 - 28 = \boxed{026}$

 





\textbf{Solution 7}

Notice that $50=5^2+5^2=7^2+1^2$, and by the size of the diagram, it seems reasonable that $OA$ represents $5^2+5^2$, and $OC$ means the $7^1+1^2$, and indeed, the values work ($7-5=2$ and $5+1=6$), so $OB^2=5^2+1^2=\boxed{026}$

 Note: THIS IS NOT A RELIABLE SOLUTION, as diagrams on the tests are not usually drawn to scale.

 




	\item (\textit{AIME I 1983}) Suppose that the sum of the squares of two complex numbers $x$ and $y$ is $7$ and the sum of the cubes is $10$. What is the largest real value that $x + y$ can have?


 
 \textbf{Answer: }004



\textbf{Solution 1}

One way to solve this problem is by \href{https://artofproblemsolving.com/wiki/index.php?title=Substitution}{substitution}. We have

 $x^2+y^2=(x+y)^2-2xy=7$ and$x^3+y^3=(x+y)(x^2-xy+y^2)=(7-xy)(x+y)=10$

 Hence observe that we can write $w=x+y$ and $z=xy$.

 This reduces the equations to $w^2-2z=7$ and$w(7-z)=10$.

 Because we want the largest possible $w$, let's find an expression for $z$ in terms of $w$.

 $w^2-7=2z \implies z=\frac{w^2-7}{2}$.

 Substituting, $w^3-21w+20=0$, which factorizes as $(w-1)(w+5)(w-4)=0$  (the \href{https://artofproblemsolving.com/wiki/index.php?title=Rational\_Root\_Theorem}{Rational Root Theorem} may be used here, along with synthetic division).

 The largest possible solution is therefore $x+y=w=\boxed{004}$.

 





\textbf{Solution 2}

An alternate way to solve this is to let $x=a+bi$ and $y=c+di$.

 Because we are looking for a value of $x+y$ that is real, we know that $d=-b$, and thus $y=c-bi$.

 Expanding $x^2+y^2=7+0i$ will give two equations, since the real and imaginary parts must match up.

 $(a+bi)^2+(c-bi)^2=7+0i$

 $(a^2+c^2-2b^2)+(2ab-2cb)i=7+0i$

 Looking at the imaginary part of that equation, $2ab-2cb=0$, so $a=c$, and $x$ and $y$ are actually complex conjugates.

 Looking at the real part of the equation and plugging in $a=c$, $2a^2-2b^2=7$, or $2b^2=2a^2-7$.

 Now, evaluating the real part of $(a+bi)^3+(a-bi)^3$, which equals $10$ (ignoring the odd powers of $i$, since they would not result in something in the form of $10+0i$):

 $a^3+3a(bi)^2+a^3+3a(-bi)^2=10$

 $2a^3-6ab^2=10$

 Since we know that $2b^2=2a^2-7$, it can be plugged in for $b^2$ in the above equation to yield:

 $2a^3-3a(2a^2-7)=10$

 $-4a^3+21a=10$

 $4a^3-21a+10=0$

 Since the problem is looking for $x+y=2a$ to be a positive integer, only positive half-integers (and whole-integers) need to be tested.  From the Rational Roots theorem, $a=10, a=5, a=\frac{5}{2}$ all fail, but $a=2$ does work.  Thus, the real part of both numbers is $2$, and their sum is $\boxed{004}$.

 





\textbf{Solution 3}

Begin by assuming that $x$ and $y$ are roots of some polynomial of the form $w^2+bw+c$, such that by \href{https://artofproblemsolving.com/wiki/index.php?title=Vieta\%27s\_Formulas}{Vieta's Formulas} and some algebra (left as an exercise to the reader), $b^2-2c=7$ and $3bc-b^3=10$.Substituting $c=\frac{b^2-7}{2}$, we deduce that $b^3-21b-20=0$, whose roots are $-4$, $-1$, and $5$.Since $-b$ is the sum of the roots and is maximized when $b=-4$, the answer is $-(-4)=\boxed{004}$.

 





\textbf{Solution 4}

$x^3 + y^3 = 10 = (x+y)(x^2-xy+y^2) = (x+y)(7-xy) \implies xy = 7 - \frac{10}{x+y}.$  Also, $(x+y)^3 = x^3 + 3x^2y+3xy^2+y^3 = 10 + 3xy(x+y).$  Substituting our above into this, we get $10 + 3(7-\frac{10}{x+y})(x+y) = 21x+21y-20 = (x+y)^3$.  Letting $p = x+y$, we have that $p^3 - 21p + 20 = 0$.  Testing $p = 1$, we find that this is a root, to get $(p-1)(p^2+p-20) = 0 \implies p = -5, 1, 4 \implies \boxed{4}$

 




	\item (\textit{AIME I 1983}) Let $a_n=6^{n}+8^{n}$. Determine the remainder on dividing $a_{83}$ by $49$.


 
 \textbf{Answer: }035



\textbf{Solution 1}

Firstly, we try to find a relationship between the numbers we're provided with and $49$. We notice that $49=7^2$, and both $6$ and $8$ are greater or less than $7$ by $1$. 

 Thus, expressing the numbers in terms of $7$, we get $a_{83} = (7-1)^{83}+(7+1)^{83}$.

 Applying the \href{https://artofproblemsolving.com/wiki/index.php?title=Binomial\_Theorem}{Binomial Theorem}, half of our terms cancel out and we are left with $2(7^{83}+3403\cdot7^{81}+\cdots + 83\cdot7)$. We realize that all of these terms are divisible by $49$ except the final term.

 After some quick division, our answer is $\boxed{035}$.

 \textbf{Clearer explanation of Solution 1} The key idea is to notice that $49=7^2$, and we need a way to subtract large multiples of $49$ from both $6^{83}$ and $8^{83}$ to make these numbers smaller.

 
Write $6$ as $7-1$ and $8$ as $7+1$, then do binomial expansion. In both $(7-1)^{83}$ and $(7+1)^{83}$, all the terms containing a greater power of $7$ than $7^2$ can be neglected because they are all multiples of $49$. Only the last two terms need to be considered.

 
The last two terms for $(7-1)^{83}$ is $83\times{7}\times{(-1)^{82}}+(-1)^{83}$.

 The last two terms for $(7+1)^{83}$ is $83\times{7}\times{(1)^{82}}+(1)^{83}$.

 
Add them together. $83\times{7}\times{(-1)^{82}}+(-1)^{83}+83\times{7}\times{(1)^{82}}+(1)^{83}$

 $=2\times 83\times 7$

 $=1162$

 
The remainder of $1162$ upon division by $49$ is $35$.

 





\textbf{Solution 2}

Since $\phi(49) = 42$ (see \href{https://artofproblemsolving.com/wiki/index.php?title=Euler\%27s\_totient\_function}{Euler's totient function}), \href{https://artofproblemsolving.com/wiki/index.php?title=Euler\%27s\_Totient\_Theorem}{Euler's Totient Theorem} tells us that $a^{42} \equiv 1 \pmod{49}$ where $\text{gcd}(a,49) = 1$. Thus $6^{83} + 8^{83} \equiv 6^{2(42)-1}+8^{2(42)-1}$ $\equiv 6^{-1} + 8^{-1} \equiv \frac{8+6}{48}$ $\equiv \frac{14}{-1}\equiv \boxed{035} \pmod{49}$.

 
                \begin{itemize}
                    \item Alternatively, we could have noted that $a^b\equiv a^{b\pmod{\phi{(n)}}}\pmod n$. This way, we have $6^{83}\equiv 6^{83\pmod {42}}\equiv 6^{-1}\pmod {49}$, and can finish the same way.
                \end{itemize}


                





\textbf{Solution 3}

$6^{83} + 8^{83} = (6+8)(6^{82}-6^{81}8+\ldots-8^{81}6+8^{82})$

 Because $7|(6+8)$, we only consider  $6^{82}-6^{81}8+\ldots-8^{81}6+8^{82}  \pmod{7}$

 $6^{82}-6^{81}8+\ldots-8^{81}6+8^{82} \equiv (-1)^{82} - (-1)^{81}+ \ldots - (-1)^1 + 1 = 83 \equiv 6 \pmod{7}$

 $6^{83} + 8^{83} \equiv 14 \cdot 6 \equiv \boxed{035} \pmod{49}$

 


 





\textbf{Solution 4 last resort (bash)}

Repeat the steps of taking modulo $49$ after reducing the exponents over and over again until you get a residue of $49,$ namely $35.$ This bashing takes a lot of time but it isn’t too bad. ~peelybonehead

 




	\item (\textit{AIME I 1983}) Twenty five of King Arthur's knights are seated at their customary round table. Three of them are chosen - all choices being equally likely - and are sent off to slay a troublesome dragon. Let $P$ be the probability that at least two of the three had been sitting next to each other. If $P$ is written as a fraction in lowest terms, what is the sum of the numerator and denominator?


 
 \textbf{Answer: }057



\textbf{Solution 1}

We can use \href{https://artofproblemsolving.com/wiki/index.php?title=Complementary\_counting}{complementary counting}, by finding the probability that none of the three knights are sitting next to each other and subtracting it from $1$.

 Imagine that the $22$ other (indistinguishable) people are already seated, and fixed into place. 

 We will place $A$, $B$, and $C$ with and without the restriction.

 There are $22$ places to put $A$, followed by $21$ places to put $B$, and $20$ places to put $C$ after $A$ and $B$. Hence, there are $22\cdot21\cdot20$ ways to place $A, B, C$ in between these people with restrictions.

 Without restrictions, there are $22$ places to put $A$, followed by $23$ places to put $B$, and $24$ places to put $C$ after $A$ and $B$. Hence, there are $22\cdot23\cdot24$ ways to place $A,B,C$ in between these people without restrictions.

 Thus, the desired probability is $1-\frac{22\cdot21\cdot20}{22\cdot23\cdot24}=1-\frac{420}{552}=1-\frac{35}{46}=\frac{11}{46}$, and the answer is $11+46=\boxed{057}$.

 





\textbf{Solution 2}

There are $(25-1)! = 24!$ configurations for the knights about the table (since configurations that are derived from each other simply by a rotation are really the same, and should not be counted multiple times). 

 There are ${3\choose 2} = 3$ ways to pick a pair of knights from the trio, and there are $2! = 2$ ways to determine in which order they are seated. Since these two knights must be together, we let them be a single entity, so there are $(24-1)! = 23!$ configurations for the entities. 

 However, this \href{https://artofproblemsolving.com/wiki/index.php?title=Principle\_of\_Inclusion-Exclusion}{overcounts} the instances in which the trio sits together; when all three knights sit together, then two of the pairs from the previous case are counted. However, we only want to count this as one case, so we need to subtract the number of instances in which the trio sits together (as a single entity). There are $3! = 6$ ways to determine their order, and there are $(23-1)! = 22!$ configurations.

 Thus, the required probability is $\frac{2 \times 3 \times 23! - 6 \times 22!}{24!} = \frac{11}{46}$, and the answer is $\boxed{057}$.

 





\textbf{Solution 3}

Number the knights around the table from $1$ to $25$. The total number of ways to pick the knights is 
 \[{25\choose 3} = \frac{25\cdot24\cdot23}{3\cdot2\cdot1} = 25\cdot23\cdot4\]There are two possibilities: either all three sit next to each other, or two sit next to each other and one is not sitting next to the other two.

 Case 1:  All three sit next to each other.  In this case, you are picking $(1,2,3)$, $(2,3,4)$, $(3,4,5)$, $(4,5,6)$, ...,$(23,24,25)$, $(24,25,1)$, $(25,1,2)$.  This makes $25$ combinations.

 Case 2:  Like above, there are $25$ ways to pick the pair of knights sitting next to each other.  Once a pair is picked, you cannot pick either of the two adjacent knights. (For example, if you pick $(5,6)$, you may not pick 4 or 7.) Thus, there are $25-4=21$ ways to pick the third knight, for a total of $25\cdot21$ combinations.

 Thus, you have a total of $25 + (25\cdot21) = 25\cdot22$ allowable ways to pick the knights.    

 The probability is $\frac{25\cdot22}{25\cdot23\cdot4} = \frac{11}{46}$, and the answer is $\boxed{057}$.

 





\textbf{Solution 4}

Pick an arbitrary spot for the first knight. Then pick spots for the next two knights in order.

 Case 1: The second knight sits next to the first knight. There are $2$ possible places for this out of $24$, so the probability of this is $\frac{1}{12}$. We do not need to consider the third knight.

 Case 2: The second knight sits two spaces apart from the first knight. There are $2$ possible places for this out of $24$, so the probability is $\frac{1}{12}$. Then there are $3$ places out of a remaining $23$ for the third knight to sit, so the total probability for this case is $\frac{1}{12} \times \frac{3}{23}$.

 Case 3: The second knight sits three or more spaces apart from the first knight. There are $20$ possible places for this out of $24$, so the probability is $\frac{5}{6}$. Then there are $4$ places to put the last knight out of $23$, so the total probability for this case is $\frac{5}{6}\times\frac{4}{23}$.

 Now we add the probabilities to get the total: 

 
 \[\frac{1}{12}+\frac{1}{12}\times\frac{3}{23}+\frac{5}{6}\times\frac{4}{23}=\frac{1}{12}\times\frac{1}{23}\left(23+3+40\right)=\frac{66}{12\times 23}\]
 \[=\frac{6\times 11}{6\times 2 \times 23}=\frac{11}{46}\] so the answer is $\boxed{057}$.

 





\textbf{Solution 5}

We simplify this problem by using complementary counting and fixing one knight in place. Then, either a knight can sit two spaces apart from the fixed knight, or a knight can sit more than two spaces apart from the fixed knight. The probability is then $\frac{24\left(23\right)-\left[2\left(20\right)+20\left(19\right)\right]}{24\left(23\right)}=\frac{11}{46}$, so the answer is $11+46=\boxed{057}$.

 





\textbf{Solution 6 (stars and bars)}

Let $K_1, K_2, K_3$ be the knights in clockwise order. Let $A$ be the distance between $K_1$ and $K_2$, $B$ the distance between $K_2$ and $K_3$ and $C$ the distance between $K_3$ and $K_1$. $A + B + C = 25$ and $A, B, C \geq 1$. In order to use stars and bars the numbers must be greater than or equal to 0 instead of 1, so we define $A_1 = A - 1, B_1 = B - 1, C_1 = C - 1$. $A_1 + B_1 + C_1 = 22$, so by stars and bars there are $\binom{22 + 3 - 1}{3 - 1} = 276$ possibilities.

 The condition is not satisfied if $A, B, C \geq 2$, so we can use complementary counting. Let $A_2 = A - 2, B_2 = B - 2, C_2 = C - 2$. $A_2 + B_2 + C_2 = 19$, and by stars and bars there are $\binom{19 + 3 - 1}{3 - 1} = 210$ possibilities. This means there are $276 - 210 = 66$ possibilities where the condition is satisfied, so the probability is $\frac{11}{46}$, resulting in $\boxed{057}$.

 





\textbf{Solution 7}

There are $\binom{25}{3}=\frac{25 \cdot 24 \cdot 23}{3 \cdot 2}=2300$ ways to chose $3$ knights out of $25$ knights.

 To ensure at least $2$ adjacent knights are chosen, first choose $1$ of the $25$ pairs of adjacent knights. After choosing the adjacent pair of knights, there are $23$ ways to choose the third knight. There are $25$ choices of $3$ knights sitting consecutively, which are counted twice. For example, choosing $(1,2)$ first, then choosing $3$ is the same as choosing $(2,3)$ first, then choosing $1$. Therefore, there are $25 \cdot 23 -25=550$ ways to choose $3$ knights where at least $2$ of the $3$ had been sitting next to each other.
 \[P=\frac{550}{2300}=\frac{11}{46}\]The answer is $11+46=\boxed{\textbf{057}}$

 





\textbf{Solution 8}

There are two possible scenarios satisfying this condition, one where two knights sit next to each other, while the third knight is lonely, and the other is where all three knights sit next to each other. On the numerator will be the total number of ways to position the knights, and on the denominator will be the number of ways to choose three knights.

 Denominator: $\binom{25}{3}$ ways to choose three knights

 Numerator:

 $\implies$ $Case_1:$ Two knights are next to each other, third is lonely :(

 $25$ ways to situate two knights, $21$ ways for third knight

 
$\implies$ $Case_2:$ All three knights are next to each other

 $25$ ways

 So, we have $P = \frac{21 \cdot 25+25}{\binom{25}{3}}$$= \frac{22 \cdot 25 \cdot 6}{25 \cdot 24 \cdot 23}$$= \frac{11}{46}$, and our answer is $11+46=\boxed{057}$

 





\textbf{Solution 9}

Lets use complementary counting. First, we must find the probability that no $3$ chosen people are sitting next to each other. We claim that if there are $n$ people around the circle, and $k$ chosen people, this probability is simply 
 \[\dfrac{n}{k}\dbinom{n-k-1}{k-1}\] Now, we can just plug it in to get a total of $1750$ ways. The denominator is simply $\dbinom{25}{3}=2300$. Now, our probability is simply $\dfrac{1750}{2300}$. But, this is the opposite of what we want. We want $1-\dfrac{1750}{2300} = \dfrac{11}{46} \Longrightarrow \boxed{057}$.

 -jb2015007

 




	\item (\textit{AIME I 1983}) What is the largest $2$-digit prime factor of the integer $n = {200\choose 100}$?


 
 \textbf{Answer: }061



\textbf{Solution 1}

Expanding the \href{https://artofproblemsolving.com/wiki/index.php?title=Combination}{binomial coefficient}, we get ${200 \choose 100}=\frac{200!}{100!100!}$. Let the required prime be $p$; then $10 \le p < 100$. If $p > 50$, then the factor of $p$ appears twice in the denominator. Thus, we need $p$ to appear as a factor at least three times in the numerator, so $3p<200$. The largest such prime is $\boxed{061}$, which is our answer.

 





\textbf{Solution 1.5, half (desperate)}

The only way a single prime number will be left out after all the cancelation will have to satisfy the condition that out of all the multiples of the prime number we want to find, 2 of the multiples will have to be between 100 and 200, $61$ is good for this solution (if you know your primes) because $61 \cdot 2$ and $61 \cdot 3$ both lie in the interval $(100, 200)$.

 The reason for this is because if 2 copies of the prime $p$ we want to find are on the numerator and one copy on the denominator (AKA it contains $p^2$ and also $\dfrac{1}{p}$), there will only be 1 copy of the prime we want to find when all the cancellation is made when reducing the original fraction (only $p$ is left after cancellation).

 





\textbf{Solution 2: Full Thinking Process of Solution 1.5}

We want the largest two‑digit prime factor of
 \[{200\choose100}=\frac{200\cdot199\cdot\ldots\cdot101}{100\cdot99\cdot\ldots\cdot1}.\]

 Let $p$ be a two‑digit prime factor of this number. Any such $p$ that remains after cancellation must appear at least twice in the numerator, because one factor of $p$ in a multiple of $p$ in the numerator will cancel with itself in the denominator. Therefore, there must be at least two multiples of $p$ between $101$ and $200$.

 Write such a multiple as $kp$, where $k$ is a positive integer. Then
 \[101\le kp\le200.\]

 Since $p$ is two‑digit, $10\le p\le99$. Dividing the inequality by $k$ gives
 \[\frac{101}{k}\le p\le\frac{200}{k}.\]

 For $k=1$, $\frac{101}{1}\le p\le\frac{200}{1}$, which gives $101\le p\le200$, impossible because $p$ is two‑digit.

 For $k=2$, $\frac{101}{2}\le p\le\frac{200}{2}$, which gives $50.5\le p\le100$, so
 \[51\le p\le99.\]

 For $k=3$, $\frac{101}{3}\le p\le\frac{200}{3}$, which gives $33.67\le p\le66.67$, so
 \[34\le p\le66.\]

 Notice that for larger $k$, the upper bound is smaller, so larger $k$ minimize the maximum possible value of $p$.

 Thus the possible range of $p$ that has at least two multiples in $[101,200]$ is the intersection
 \[[51,99]\cap[34,66]=[51,66].\]

 The largest prime in $[51,66]$ is $61$.

 Therefore the largest two‑digit prime factor of ${200\choose100}$ is $\boxed{61}$.

 Note: If $p$ was less than or equal to 50, then for $k=2$, $kp$ would not be in the range $[101,200]$. But since we have found a prime greater than 50, $k=2$ and $k=3$ works.

 





\textbf{Solution 3: Clarification of Solution 1}

We know that 
 \[{200\choose100}=\frac{200!}{100!100!}\]Since $p<100$, there is at least $1$ factor of $p$ in each of the $100!$ in the denominator. Thus there must be at least $3$ factors of $p$ in the numerator $200!$ for $p$ to be a factor of $n=\frac{200!}{100!100!}$. (Note that here we assume the minimum because as $p$ goes larger in value, the number of factors of $p$ in a number decreases,)

 So basically, $p$ is the largest prime number such that
 \[\left \lfloor\frac{200}{p}\right \rfloor>3\]Since $p<\frac{200}{3}=66.66...$, the largest prime value for $p$ is $p=\boxed{61}$

 \textbf{Note} Similar to 2023 MATHCOUNTS State Sprint \#25

 \textbf{Comment} which came first, 1985 or 2023?

 




	\item (\textit{AIME I 1983}) Find the minimum value of $\frac{9x^2\sin^2 x + 4}{x\sin x}$ for $0 < x < \pi$.


 
 \textbf{Answer: }012



\textbf{Solution 1}

Let $y=x\sin{x}$. We can rewrite the expression as $\frac{9y^2+4}{y}=9y+\frac{4}{y}$.

 Since $x>0$, and $\sin{x}>0$ because $0< x<\pi$, we have $y>0$. So we can apply \href{https://artofproblemsolving.com/wiki/index.php?title=AM-GM}{AM-GM}:

 
 \[9y+\frac{4}{y}\ge 2\sqrt{9y\cdot\frac{4}{y}}=12\]

 The equality holds when $9y=\frac{4}{y}\Longleftrightarrow y^2=\frac49\Longleftrightarrow y=\frac23$.

 Therefore, the minimum value is $\boxed{012}$. This is reached when we have $x \sin{x} = \frac{2}{3}$ in the original equation (since $x\sin x$ is continuous and increasing on the interval $0 \le x \le \frac{\pi}{2}$, and its range on that interval is from $0 \le x\sin x \le \frac{\pi}{2}$, this value of $\frac{2}{3}$ is attainable by the \href{https://artofproblemsolving.com/wiki/index.php?title=Intermediate\_Value\_Theorem}{Intermediate Value Theorem}).

 





\textbf{Solution 2}

We can rewrite the numerator to be a perfect square by adding $-\dfrac{12x \sin x}{x \sin x}$. Thus, we must also add back $12$.

 This results in $\dfrac{(3x \sin x-2)^2}{x \sin x}+12$.

 Thus, if $3x \sin x-2=0$, then the minimum is obviously $12$. We show this possible with the same methods in Solution 1; thus the answer is $\boxed{012}$.

 





\textbf{Solution 3 (uses calculus)}

Let $y = x\sin{x}$ and rewrite the expression as $f(y) = 9y + \frac{4}{y}$, similar to the previous solution. To minimize $f(y)$, take the \href{https://artofproblemsolving.com/wiki/index.php?title=Derivative}{derivative} of $f(y)$ and set it equal to zero. 

 The derivative of $f(y)$, using the Power Rule, is

 $f'(y)$ = $9 - 4y^{-2}$

 $f'(y)$ is zero only when $y = \frac{2}{3}$ or $y = -\frac{2}{3}$. It can further be verified that $\frac{2}{3}$ and $-\frac{2}{3}$ are relative minima by finding the derivatives at other points near the critical points, or by checking that the second derivative $f''(y)=8y^{-3}$ is positive. However, since $x \sin{x}$ is always positive in the given domain, $y = \frac{2}{3}$. Therefore, $x\sin{x}$ = $\frac{2}{3}$, and the answer is $\frac{(9)(\frac{2}{3})^2 + 4}{\frac{2}{3}} = \boxed{012}$.

 





\textbf{Solution 4 (also uses calculus)}

As above, let $y = x\sin{x}$. Add $\frac{12y}{y}$ to the expression and subtract $12$, giving $f(x) = \frac{(3y+2)^2}{y} - 12$. Taking the \href{https://artofproblemsolving.com/wiki/index.php?title=Derivative}{derivative} of $f(x)$ using the \href{https://artofproblemsolving.com/wiki/index.php?title=Chain\_Rule}{Chain Rule} and \href{https://artofproblemsolving.com/wiki/index.php?title=Quotient\_Rule}{Quotient Rule}, we have $\frac{\text{d}f(x)}{\text{d}x} = \frac{6y(3y+2)-(3y+2)^2}{y^2}$. We find the minimum value by setting this to $0$. Simplifying, we have $6y(3y+2) = (3y+2)^2$ and $y = \pm{\frac{2}{3}} = x\sin{x}$. Since both $x$ and $\sin{x}$ are positive on the given interval, we can ignore the negative root. Plugging $y = \frac{2}{3}$ into our expression for $f(x)$, we have $\frac{(3(\frac{2}{3})+2)^2}{y}-12 = \frac{16}{\left(\frac{2}{3}\right)}-12 = \boxed{012}$.

 





\textbf{Solution 5}

Set $\frac{9x^2\sin^2 x + 4}{x\sin x}$ equal to $y$. Then multiply by $x\sin x$ on both sides to get $9x^2\sin^2 x + 4 = y\cdot x\sin x$.  We then subtract $yx\sin x$ from both sides to get $9x^2\sin^2 x + 4 - yx\sin x = 0$.  This looks like a quadratic so set $z= x\sin x$ and use quadratic equation on $9z^2 - yz + 4 = 0$ to see that $z = \frac{y\pm\sqrt{y^2-144}}{18}$. We know that $y$ must be an integer and as small as it can be, so $y$ = 12. We plug this back in to see that $x\sin x = \frac{2}{3}$ which we can prove works using methods from solution 1. This makes the answer $\boxed{012}$

 
-awesomediabrine

 





\textbf{Solution 6}

Seeing that we need to minimize, we think inequalities, and seeing squares, we think \href{https://artofproblemsolving.com/wiki/index.php?title=RMS-AM-GM-HM}{RMS-AM-GM-HM}. From this inequality, we know that $\sqrt{\frac{(3x\sin x)^2+2^2}{2}} \geq \sqrt{(3x\sin x)(2)}$, with equality holding when $3x\sin x=2$. From this inequality, we can see the following:\begin{align*}\sqrt{\frac{(3x\sin x)^2+2^2}{2}} \geq \sqrt{(3x\sin x)(2)} \\\frac{9x^2\sin^2x+4}{2} \geq 6x\sin x \\\frac{9x^2\sin^2x+4}{x\sin x} \geq 12\end{align*}
We can prove that the equality condition is possible as in Solution $1$. Thus, our answer is $\boxed{012}$.

 




	\item (\textit{AIME I 1983}) The numbers $1447$, $1005$ and $1231$ have something in common: each is a $4$-digit number beginning with $1$ that has exactly two identical digits. How many such numbers are there?


 
 \textbf{Answer: }432



\textbf{Solution 1}

Suppose that the two identical digits are both $1$. Since the thousands digit must be $1$, only one of the other three digits can be $1$. This means the possible forms for the number are

 $11xy,\qquad 1x1y,\qquad1xy1$Because the number must have exactly two identical digits, $x\neq y$, $x\neq1$, and $y\neq1$. Hence, there are $3\cdot9\cdot8=216$ numbers of this form.

 Now suppose that the two identical digits are not $1$. Reasoning similarly to before, we have the following possibilities:

 $1xxy,\qquad1xyx,\qquad1yxx.$Again, $x\neq y$, $x\neq 1$, and $y\neq 1$. There are $3\cdot9\cdot8=216$ numbers of this form.

 Thus the answer is $216+216=\boxed{432}$.

 





\textbf{Solution 2}

Consider all four-digit sequences with exactly two identical digits. By symmetry, $\frac{1}{10}$ of these sequences start with $1$.

 Let $A$ be the digit that appears twice, and $B > C$ be the digits which appear once. Then there are $10$ choices for $A$ and $\binom{9}{2} = 36$ choices for $(B, C)$.  Therefore, there are $10 \cdot 36 = 360$ choices for $(A, B, C)$.

 For each choice of $(A, B, C)$, there are $\frac{4!}{2!} = 12$ permutations of $AABC$.  Therefore, there are $12 \cdot 360 = 4320$ four-digit sequences with exactly two identical digits.  Of these sequences, $\frac{1}{10} \cdot 4320 = \boxed{432}$ begin with $1$.

 





\textbf{Solution 3 (Complementary Counting)}

We'll use complementary counting. We will split up into $3$ cases: (1) no number is repeated, (2) $2$ numbers are repeated, and $2$ other numbers are repeated, (3) $3$ numbers are repeated, or (4) $4$ numbers are repeated.

 Case 1:There are $9$ choices for the hundreds digit (it cannot be $1$), $8$ choices for the tens digit (it cannot be $1$ or what was chosen for the hundreds digit), and $7$ for the units digit. This is a total of $9\cdot8\cdot7=504$ numbers.

 Case 2:One of the three numbers must be $1$, and the other two numbers must be the same number, but cannot be $1$ (That will be dealt with in case 4). There are $3$ choices to put the $1$, and there are $9$ choices (any number except for $1$) to pick the other number that is repeated, so a total of $3\cdot9=27$ numbers.

 Case 3: We will split it into $2$ subcases: one where $1$ is repeated $3$ times, and one where another number is repeated $3$ times.When $1$ is repeated $3$ times, then one of the digits is not $1$. There are $9$ choices for that number, and $3$ choices for its location,so a total of $9\cdot3=27$ numbers.When a number other than $1$ is repeated $3$ times, then there are $9$ choices for the number, and you don't have any choices on where to put that number. So in Case 3 there are $27+9=36$ numbers

 Case 4:There is only $1$ number: $1111$.

 There are a total of $1000$ $4$-digit numbers that begin with $1$ (from $1000$ to $1999$), so by complementary counting you get $1000-(504+27+36+1)=\boxed{432}$ numbers.

 Solution by Kinglogic

 





\textbf{Solution 4 (Easy casework)}

Let us proceed by casework. 

 Case 1:We will count the amount of numbers that have two identical digits that are not one. The thousands digit is fixed, and we are choosing two spots to hold two identical digits that are chosen from $0, 2-9$, which is $9$ options. For the last digit, their are $8$ possibilities since it can be neither $1$ or the other number that is chosen. The final outcome is ${3}\choose{2}$ $* 9 * 8 = 216$ possibilities for this case.

 Case 2:The last case will be the amount of numbers that have two identical digits thare are $1$. There are ${3}\choose{1}$ places to pick the $1$. For the other $2$ digits, there are $9$ and $8$ options respectively. Thus, we have $3 * 9 * 8 = 216$.Summing the two cases, we get $216 + 216 = \boxed{432}$.

 




	\item (\textit{AIME I 1983}) The solid shown has a square base of side length $s$. The upper edge is parallel to the base and has length $2s$. All other edges have length $s$. Given that $s=6\sqrt{2}$, what is the volume of the solid?


 
                    \begin{figure}[H]
                        \centering
                        \includegraphics[width=0.5\textwidth]{Problem/AIME_1983_I_11_Problem_Figure_1.png}
                    \end{figure}
                    


 
 \textbf{Answer: }288



\textbf{Solution 1}

First, we find the height of the solid by dropping a perpendicular from the midpoint of $AD$ to $EF$. The hypotenuse of the triangle formed is the \href{https://artofproblemsolving.com/wiki/index.php?title=Median}{median} of equilateral triangle $ADE$, and one of the legs is $3\sqrt{2}$. We apply the Pythagorean Theorem to deduce that the height is $6$.

 
                    \begin{figure}[H]
                        \centering
                        \includegraphics[width=0.5\textwidth]{Solution/AIME_1983_I_11_Solution_Figure_1.png}
                    \end{figure}
                    Next, we complete the figure into a triangular prism, and find its volume, which is $\frac{6\sqrt{2}\cdot 12\sqrt{2}\cdot 6}{2}=432$.

 Now, we subtract off the two extra \href{https://artofproblemsolving.com/wiki/index.php?title=Pyramid}{pyramids} that we included, whose combined volume is $2\cdot \left( \frac{6\sqrt{2}\cdot 3\sqrt{2} \cdot 6}{3} \right)=144$.

 Thus, our answer is $432-144=\boxed{288}$.

 





\textbf{Solution 2}


                    \begin{figure}[H]
                        \centering
                        \includegraphics[width=0.5\textwidth]{Solution/AIME_1983_I_11_Solution_Figure_1.png}
                    \end{figure}
                    Extend $EA$ and $FB$ to meet at $G$, and $ED$ and $FC$ to meet at $H$. Now, we have a regular tetrahedron $EFGH$, which by symmetry has twice the volume of our original solid. This tetrahedron has side length $2s = 12\sqrt{2}$. Using the formula for the volume of a regular tetrahedron, which is $V = \frac{\sqrt{2}S^3}{12}$, where S is the side length of the tetrahedron, the volume of our original solid is:

 $V = \frac{1}{2} \cdot \frac{\sqrt{2} \cdot (12\sqrt{2})^3}{12} = \boxed{288}$.

 





\textbf{Solution 3}

We can also find the volume by considering horizontal cross-sections of the solid and using calculus. As in Solution 1, we can find that the height of the solid is $6$; thus, we will integrate with respect to height from $0$ to $6$, noting that each cross section of height $dh$ is a rectangle. The volume is then $\int_0^h(wl) \ \text{d}h$, where $w$ is the width of the rectangle and $l$ is the length. We can express $w$ in terms of $h$ as $w=6\sqrt{2}-\sqrt{2}h$ since it decreases linearly with respect to $h$, and $l=6\sqrt{2}+\sqrt{2}h$ since it similarly increases linearly with respect to $h$. Now we solve:
 \[\int_0^6(6\sqrt{2}-\sqrt{2}h)(6\sqrt{2}+\sqrt{2}h)\ \text{d}h =\int_0^6(72-2h^2)\ \text{d}h=72(6)-2\left(\frac{1}{3}\right)\left(6^3\right)=\boxed{288}\].

 





\textbf{Solution 4}

Draw an altitude from a vertex of the square base to the top edge. By using $30,60, 90$ triangle ratios, we obtain that the altitude has a length of $3 \sqrt{6}$, and that little portion that hangs out has a length of $3\sqrt2$. This is a triangular pyramid with a base of $3\sqrt6, 3\sqrt6, 3\sqrt2$, and a height of $3\sqrt{2}$. Since there are two of these, we can compute the sum of the volumes of these two to be $72$. Now we are left with a triangular prism with a base of dimensions $3\sqrt6, 3\sqrt6, 3\sqrt2$ and a height of $6\sqrt2$. We can compute the volume of this to be 216, and thus our answer is $\boxed{288}$.

 pi\_is\_3.141

 





\textbf{Solution 5}

From solution 1, the height of the solid is $6$. Construct a triangular prism with base ABCD, and with the height of the solid. The volume of this triangular prism is $( 6\sqrt{2} )^2 \cdot 6 \cdot 1/2 = 216$ Notice that the solid is symetrical, and if you remove the triangular prism described earlier and combine the two halves on each side of it, you will get a regular tetrahedron. This tetrahedron has sidelength of $6\sqrt{2}$ and since the formula for the volume of a regular tetrahedron is $\frac{s^3}{6\sqrt{2}}$ the volume of this tetrahedron is $72$. $216 + 72 = \boxed{288}$. 

 




	\item (\textit{AIME I 1983}) Diameter $AB$ of a circle has length a $2$-digit integer (base ten). Reversing the digits gives the length of the perpendicular chord $CD$. The distance from their intersection point $H$ to the center $O$ is a positive rational number. Determine the length of $AB$.


 
                    \begin{figure}[H]
                        \centering
                        \includegraphics[width=0.5\textwidth]{Problem/AIME_1983_I_12_Problem_Figure_1.png}
                    \end{figure}
                    


 
 \textbf{Answer: }065



\textbf{Solution}

Let $AB=10x+y$ and $CD=10y+x$. It follows that $CO=\frac{AB}{2}=\frac{10x+y}{2}$ and $CH=\frac{CD}{2}=\frac{10y+x}{2}$. Scale up this triangle by 2 to ease the arithmetic. Applying the \href{https://artofproblemsolving.com/wiki/index.php?title=Pythagorean\_Theorem}{Pythagorean Theorem} on $2CO$, $2OH$ and $2CH$, we deduce
 \[(2OH)^2=\left(10x+y\right)^2-\left(10y+x\right)^2=99(x+y)(x-y)\] 

 Because $OH$ is a positive rational number and $x$ and $y$ are integral, the quantity $99(x+y)(x-y)$ must be a perfect square. Hence either $x-y$ or $x+y$ must be a multiple of $11$, but as $x$ and $y$ are different digits, $1+0=1 \leq x+y \leq 9+9=18$, so the only possible multiple of $11$ is $11$ itself. However, $x-y$ cannot be 11, because both must be digits. Therefore, $x+y$ must equal $11$ and $x-y$ must be a perfect square. The only pair $(x,y)$ that satisfies this condition is $(6,5)$, so our answer is $\boxed{065}$. (Therefore $CD = 56$ and $OH = \frac{33}{2}$.)

 \textbf{Alternate start to solution 1} Since H is the midpoint of $CD$, by \href{https://artofproblemsolving.com/wiki/index.php?title=Power\_of\_a\_Point}{Power of a Point}, $CH^2=(AH)(BH)$. Because $AH=r-OH$ and $BH=r+OH$, where $r$ is the radius of the circle, we deduce the relation $CH^2=r^2-OH^2$. Thus $OH^2=\left(\frac{10x+y}{2}\right)^2-\left(\frac{10y+x}{2}\right)^2$. Continue as above.

 




	\item (\textit{AIME I 1983}) For $\{1, 2, 3, \ldots, n\}$ and each of its nonempty subsets a unique \textbf{alternating sum} is defined as follows. Arrange the numbers in the subset in decreasing order and then, beginning with the largest, alternately add and subtract succesive numbers. For example, the alternating sum for $\{1, 2, 3, 6,9\}$ is $9-6+3-2+1=5$ and for $\{5\}$ it is simply $5$. Find the sum of all such alternating sums for $n=7$.


 
 \textbf{Answer: }448



\textbf{Solution 1}

Let $S$ be a non-\href{https://artofproblemsolving.com/wiki/index.php?title=Empty\_set}{ empty} \href{https://artofproblemsolving.com/wiki/index.php?title=Subset}{subset} of $\{1,2,3,4,5,6\}$. 

 Then the alternating sum of $S$, plus the alternating sum of $S \cup \{7\}$, is $7$. This is because, since $7$ is the largest element, when we take an alternating sum, each number in $S$ ends up with the opposite sign of each corresponding element of $S\cup \{7\}$.

 Because there are $2^{6}=64$ of these pairs of sets, the sum of all possible subsets of our given set is $64 \cdot 7$, giving an answer of $\boxed{448}$.

 





\textbf{Solution 2 (almost the same as Solution 1)}

Consider a given subset $T$ of $S$ that contains $7$; then there is a subset $T'$ which contains all the elements of $T$ except for $7$, and only those elements . Since each element of $T'$ has one fewer element preceding it than it does in $T$, their signs are opposite. Thus the sum of the alternating sums of $T$ and $T'$ is equal to 7. There are $2^6$ subsets containing 7, so our answer is $7 \cdot 2^6 = \boxed{448}$.

 





\textbf{Solution 3}

Denote the desired total of all alternating sums of an $n$-element set as $S_n$. We are looking for $S_7$. Notice that all alternating sums of an $n$-element set are also alternating sums of an $n+1$-element set. However, when we go from an $n$ to $n+1$ element set, for each subset with the new element, we are adding the new element and subtracting one of the alternating sums of the $n$-element set. There are $2^n$ subsets of an $n+1$-element set that includes the new element, giving us the relationship $S_{n+1} = S_n + 2^n(n+1) - S_n = 2^n(n+1)$. When $n = 6$, we therefore get $S_ 7 = 2^6(7) = \boxed{448}$.

 





\textbf{Solution 4}

We analyze all the numbers from 1 to 7 separately to see where the number contributes its positive or negative to the sum of the alternating sums. Whenever 7 appears, which it does 64 times, it contributes a positive because it is always first. This gives a net gain of $7 \cdot 64=448$.

 If we look at when 6 appears, which it also does 64 times, whether it comes as positive or negative depends on the presence of 7. Half of the subsets with 6 have 7 resulting in subtracting 6 each time, while the other half does not have 7 adding 6 each time, so these contributions of sixes cancel each other out giving a net gain of 0.

 The same thing happens to any positive integer less than 7. This is because the determination of a positive or negative contribution is dependent on the number of larger numbers in front of it(For example, the sign of 3 is dependent on the presence of 4, 5, 6, and 7 in the subset). If the number of larger numbers is even, it gives in a positive copy while odd produces its negative. We know that the frequencies of these two cases occurring are the same because $0=(1-1)^{n}=\binom{n}{0}-\binom{n}{1}+\binom{n}{2}-...\binom{n}{n}$ via the Binomial Theorem. Therefore, all positive integers less than 7 will not have any effect and our sum will be $\boxed{448}$.

 





\textbf{Solution 5}

Let $\mathbb{N}_n := \{1, 2, 3, \dots n\}$. Let the alternating sum of a certain subset of $S$ of $\mathbb{N}_n$ be $\xi(S),$ and let 
 \[\mathcal{A}(\mathbb{N}_n) := \sum_{S \subseteq \mathbb{N}_n} \xi(S).\] We see that 
 \[\mathcal{A}(\mathbb{N}_n) = \sum_{S \subseteq \mathbb{N}_n} \xi(S) = \sum_{n \in S, S \subseteq \mathbb{N}_n} \xi(S) + \sum_{S \subseteq \mathbb{N}_{n-1}} \xi(S) = \mathcal{A}(\mathbb{N}_{n-1}) + \sum_{n \in S, S \subseteq \mathbb{N}_n} \left( n - \xi(S - \{n\}) \right),\] as if $n \in S,$ $n$ is the largest element in $S.$ Now, we know that 
 \[\sum_{n \in S, S \subseteq \mathbb{N}_n} n - \xi(S - \{n\}) = \sum_{S \subseteq \mathbb{N}_{n-1}} n - \xi(S) = n \cdot 2^{n-1} - \sum_{S \subseteq \mathbb{N}_{n-1}} \xi(S) = n \cdot 2^{n-1} - \mathcal{A}(\mathbb{N}_{n-1}),\] so 
 \[\mathcal{A}(\mathbb{N}_{n}) = n \cdot 2^{n-1}.\]  Thus, our answer (which is the $n = 7$ case) is $\mathcal{A}(\mathbb{N}_{7}) = 7 \cdot 2^6 = \boxed{448}.$

 





\textbf{Solution 6 (Lucky Guess, Not Recommended)}

Computing the first few $S_n$ gives, $S_1$ $=$ $1$, $S_2$ $=$ $4$, $S_3$ $=$ $12$, and $S_4$ $=$ $32$. From this, we think the formula for $S_n$ must be an exponential function in terms of $2^n$ because the amount of subsets that a set with $n$ elements has is $2^n$ and this is a question about subsets so we try $S_n$ = $2^n(x)$. Trying this formula for the first few $S_n$, gives away $x$ to be $n/2$ or more simply $S_n$ = $2^{n-1}(n)$ (try proving this by induction) . Substuiting $n = 7$ gives the requested answer of $S_ 7 = 2^6(7) = \boxed{448}$.

 Sidenote

 This solution is not a proof and since the AIME does not require proof, it is okay to finsh the answer like this. If you would like to see why the formula holds, please observe Solution 3.

 -ThemathCanadian

 


 




	\item (\textit{AIME I 1983}) In the adjoining figure, two circles with radii $8$ and $6$ are drawn with their centers $12$ units apart. At $P$, one of the points of intersection, a line is drawn in such a way that the chords $QP$ and $PR$ have equal length. Find the square of the length of $QP$. 


 
                    \begin{figure}[H]
                        \centering
                        \includegraphics[width=0.5\textwidth]{Problem/AIME_1983_I_14_Problem_Figure_1.png}
                    \end{figure}
                    


 
 \textbf{Answer: }130



\textbf{Solution 1}

Firstly, notice that if we reflect $R$ over $P$, we get $Q$. Since we know that $R$ is on circle $B$ and $Q$ is on circle $A$, we can reflect circle $B$ over $P$ to get another circle (centered at a new point $C$, and with radius $6$) that intersects circle $A$ at $Q$. The rest is just finding lengths, as follows.

 


 
                    \begin{figure}[H]
                        \centering
                        \includegraphics[width=0.5\textwidth]{Solution/AIME_1983_I_14_Solution_Figure_1.png}
                    \end{figure}
                    B is reflected like so
Since $P$ is the midpoint of segment $BC$, $AP$ is a median of $\triangle ABC$. Because we know $AB=12$, $BP=PC=6$, and $AP=8$, we can find the third side of the triangle using \href{https://artofproblemsolving.com/wiki/index.php?title=Stewart\%27s\_Theorem}{Stewart's Theorem} or similar approaches. We get $AC = \sqrt{56}$.

 


 
                    \begin{figure}[H]
                        \centering
                        \includegraphics[width=0.5\textwidth]{Solution/AIME_1983_I_14_Solution_Figure_1.png}
                    \end{figure}
                    The Kite is formed
Now we have a kite $AQCP$ with $AQ=AP=8$, $CQ=CP=6$, and diagonal $AC=\sqrt{56}$, and all we need is the length of the other diagonal $PQ$. The easiest way it can be found is with the Pythagorean Theorem. Let $2x$ be the length of $PQ$. Then

 
 \[\sqrt{36-x^2} + \sqrt{64-x^2} = \sqrt{56}.\]

 Solving this equation, we find that $x^2=\frac{65}{2}$, so $PQ^2 = 4x^2 = \boxed{130}.$

 





\textbf{Solution 2 (Easiest)}


                    \begin{figure}[H]
                        \centering
                        \includegraphics[width=0.5\textwidth]{Solution/AIME_1983_I_14_Solution_Figure_1.png}
                    \end{figure}
                    \textbf{ }

 Draw additional lines as indicated. Note that since triangles $AQP$ and $BPR$ are isosceles, the altitudes are also bisectors, so let $QM=MP=PN=NR=x$.

 Since $\frac{AR}{MR}=\frac{BR}{NR},$ triangles $BNR$ and $AMR$ are similar. If we let $y=BN$, we have $AM=3BN=3y$. 

 Applying the Pythagorean Theorem on triangle $BNR$, we have $x^2+y^2=36$. Similarly, for triangle $QMA$, we have $x^2+9y^2=64$.

 Subtracting, $8y^2=28\Rightarrow y^2=\frac72\Rightarrow x^2=\frac{65}2\Rightarrow QP^2=4x^2=\boxed{130}$.

 





\textbf{Solution 3 (trig bash)}

Let $QP=PR=x$. Angles $QPA$, $APB$, and $BPR$ must add up to $180^{\circ}$. By the \href{https://artofproblemsolving.com/wiki/index.php?title=Law\_of\_Cosines}{Law of Cosines}, $\angle APB=\cos^{-1}\left(\frac{{-11}}{24}\right)$. Also, angles $QPA$ and $BPR$ equal $\cos^{-1}\left(\frac{x}{16}\right)$ and $\cos^{-1}\left(\frac{x}{12}\right)$. So we have 

 $\cos^{-1}\left(\frac{x}{16}\right)+\cos^{-1}\left(\frac{{-11}}{24}\right)=180^{\circ}-\cos^{-1}\left(\frac{x}{12}\right).$ Taking the cosine of both sides, and simplifying using the addition formula for $\cos$ as well as the identity $\sin^{2}{x} + \cos^{2}{x} = 1$, gives $x^2=\boxed{130}$.

 





\textbf{Solution 4 (quickest)}

Let $QP = PR = x$. Extend the line containing the centers of the two circles to meet $R$, and to meet the other side of the large circle at a point $S$.

 The part of this line from $R$ to the point nearest to $R$ where it intersects the larger circle has length $6+(12-8)=10$.   The length of the diameter of the larger circle is $16$.

 Thus by Power of a Point in the circle passing through $Q$, $R$, and $S$, we have $x \cdot 2x = 10 \cdot (10+16) = 260$, so $x^2 = \boxed{130}$.

 


 





\textbf{Solution 5 (Pythagorean Theorem and little algebraic manipulation)}


                    \begin{figure}[H]
                        \centering
                        \includegraphics[width=0.5\textwidth]{Solution/AIME_1983_I_14_Solution_Figure_1.png}
                    \end{figure}
                    Note that the midpoint of $AB$ is $T.$ Also, since $AM,NB$ bisect $QP$ and $PR,$ respectively, $P$ is the midpoint of $MN.$ Thus, $AM+NB=2PT.$ let $AM=a,BN=b.$ This means that $a+b=2PT.$ Using Stewart's Theorem on $\Delta APB,$ the median $PT$ has length $\sqrt{14}.$ Thus, $a+b=2\sqrt{14}.$ Also, since $MP=PN$, from the Pythagorean Theorem, $8^2-a^2=6^2-b^2\implies a^2-b^2=28.$ Thus, $a-b=\frac{28}{2\sqrt{14}}=\sqrt{14}.$ We conclude that $QP=MN=\sqrt{12^2-(a-b)^2}=\sqrt{130}\implies\boxed{130}.$~pinkpig

 





\textbf{Solution 6 (Only simple geometry and algebra needed)}

Looking at Drawing 2 (by the way, we don't need point $R$), we set $AM=a$ and $BN=b$, and the desired length$QP=x=PR$. We know that a radius perpendicular to a chord bisects the chord, so $MP=\frac{x}{2}$ and $PN=\frac{x}{2}$. Draw line $AP$ and $PB$, and we see that they are radii of Circles $A$ and $B$, respectively. We can write the Pythagorean relationships $a^2+\left(\frac{x}{2}\right)^2=8^2$ for triangle $AMP$ and $b^2+\left(\frac{x}{2}\right)^2=6^2$ for triangle $BNP$. We also translate segment $MN$ down so that $N$ coincides with $B$, and form another right triangle. From that triangle, you can see that the shorter leg is on the left side, having length $a-b$, the longer leg is the same as $MN=x$, and the hypotenuse is $AB=12$. We can write the Pythagorean relationship $(a-b)^2+x^2=12^2$. Solving the system of 3 unknowns and 3 equations (One of the best ways to do it is to solve for $a$ in the first equation and $b$ in the second equation, and substitute into the third equation, get an equation only in terms of $x$, and solve), you find that $x=\sqrt{130}$, so $x^2 = \boxed{130}$.

 Solution by Kinglogic

 





\textbf{Solution 7}

The centers are collinear, you can prove it (but that is already given in the later section [Proof that R,A, and B are collinear]). Drop a perpendicular from $P$ to $AB.$ You then have 2 separate segments, separated by the foot of the altitude of $P$. Call them $a$ and $b$ respectively. Call the measure of the foot of the altitude of $P$ $h$. You then have 3 equations:

 
 \[(1)\quad a+b=12\] (this is given by the fact that the distance between the centers is 12.

 
 \[(2)\quad a^2+h^2=64\]. This is given by the fact that P is on the circle with radius 8.

 
 \[(3)\quad b^2+h^2=36\]. This is given by the fact that P is on the circle with radius 6.

 Subtract (3) from (2) to get that $a^2-b^2=28$. As per (1), then you have $a-b=\frac{7}{3}$ (4). Add (1) and (4) to get that $2a=\frac{43}{3}$. Then substitute into (1) to get $b=\frac{29}{6}$. Substitute either a or b into (2) or (3) to get $h=\sqrt{455}{6}$. Then to get $PQ=PR$ it is just $\sqrt{(b+6)^2+h^2}=\sqrt{\frac{65^2}{6^2}+\frac{455}{6^2}}=\sqrt{\frac{4680}{36}}=\sqrt{130}$.

 $PQ^2=\boxed{130}$

 -dragoon

 \textbf{Full Proof that R, A, B are collinear} 
                    \begin{figure}[H]
                        \centering
                        \includegraphics[width=0.5\textwidth]{Solution/AIME_1983_I_14_Solution_Figure_1.png}
                    \end{figure}
                    

 Let $M$ and $N$ be the feet of the perpendicular from $A$ to $PQ$ and $B$ to $PR$ respectively. It is well known that a perpendicular from the center of a circle to a chord of that circle bisects the chord, so $QM = MP = PN = NR$, since the problem told us $QP = PR$.

 We will show that $R$ lies on $AB$.

 Let $T$ be the intersection of circle centered at $B$ with $AB$. Then $BT = TA = 6$. 

 Let $P$' be the foot of the perpendicular from $T$ to $MN$. Then $TP'$ is a midline (or midsegment) in trapezoid $AMNB$, so $P'$ coincides with $P$ (they are both supposed to be the midpoint of $MN$). In other words, since $\angle TP'N = 90^\circ$, then $\angle TPN = 90^\circ$.

 Thus, $\angle TPR$ subtends a $90^\circ \times 2 = 180^\circ$ degree arc. So arc $TR$ in circle $B$ is $180^\circ$, so $TR$ is a diameter, as desired. Thus $A$, $B$, $R$ are collinear.

 NOTE: Note this collinearity only follows from the fact that $6$ is half of $12$ in the problem statement. The collinearity is untrue in general.

 





\textbf{Solution 8 (Coordinate Bash)}

We use coordinate geometry to approach this problem. Let the center of larger circle be the origin $O_1$, the smaller circle be $O_2$, and the x-axis be $O_1O_2$. Hence, we can get the two circle equations: $x^2+y^2 = 64$ and $(x-12)^2+y^2=36$.

 Let point $P$ be $(a, b)$. Noting that it lies on both circles, we can plug the coordinates into both equations:

 $a^2 +b^2 = 64$$(a-12)^2+ b^2 \Rightarrow a^2-24a+144+b^2 = 64$Substituting $a^2+b^2 = 64$ into equation 2 and solving for $a$, we get $a = \frac{43}{6}$. The problem asks us to find $QP^2$, which is congruent to $PR^2$. Using the distance formula for $P(a, b)$ and $R(18, 0)$ (by Solution 7's collinear proof), we get $PR^2 = (18-a)^2 +b^2$. Using $a^2+b^2 = 64$, we find that $b^2 = \frac{455}{36}$. Plugging the variables $a$ and $b^2$ in, we get $PR^2 = QP^2 =  \boxed{130}$~SoilMilk

 





\textbf{Solution 9 (basic solution)}

Let the center of the circle with radius $8$ be $A,$ and let the center of the one with radius $6$ be $B.$ Also, let $QP = PR = x.$ Using law of cosines on triangle $\Delta APB,$ we have that $\cos{\angle{APB}} = -\frac{{11}}{24}.$ Angle chasing gives that $\angle{QAR} = \angle{APB},$ so their cosines are the same. Applying law of cosines again on triangle $\Delta QAR,$ we have $\left(2x^2\right)=64+324-2(8)(18)\left(-\frac{11}{24}\right),$ which gives that $x^2 = \boxed{130}$

 





\textbf{Solution 10 (Spiral Sym)}

If you call X the second intersection of the two circles and the centers $O_1$, $O_2$ respectively, note that triangles $O_1 X O_2$ and $Q X R$ are similar by Spiral Sym. $P$ is the midpoint of segment $Q R$. If the midpoint of $O_1 O_2$ is $M$, then $\frac{X M}{O_1 M} = \frac{X P}{Q P}$. By Appolonius Median Length theorem, $X M = \sqrt{14}$. Note that $X P$ is simply two times the height from X to $O_1O_2$, and as a result, by Heron's formula, $X P = \frac{\sqrt{(7)(13)(5)}}{3}$, and from our ratio, $Q P = \sqrt{130}$. As a result, the square is $130$, and we are done. - sepehr2010

 





\textbf{Solution 11 (Quick)}

Let A and B denote the centers of the left and right circles respectively and $QP = PR = x.$ Observe that $\Delta AQP$ and $\Delta AQR$ share a common angle $\angle{AQP}.$ Using the law of cosines on both triangles, we get:

 
 \[(1)\quad 64=x^2+64-16x\cos{\angle{AQP}}\]
 \[(2)\quad 324=4x^2+64-32x\cos{\angle{AQP}}\]

 Subtracting $2*(1)$ from $(2),$ we get $196=2x^2-64,$ which gives $x^2=\boxed{130}$

 




	\item (\textit{AIME I 1983}) The adjoining figure shows two intersecting chords in a circle, with $B$ on minor arc $AD$. Suppose that the radius of the circle is $5$, that $BC=6$, and that $AD$ is bisected by $BC$. Suppose further that $AD$ is the only chord starting at $A$ which is bisected by $BC$. It follows that the sine of the central angle of minor arc $AB$ is a rational number. If this number is expressed as a fraction $\frac{m}{n}$ in lowest terms, what is the product $mn$?

                    \begin{figure}[H]
                        \centering
                        \includegraphics[width=0.5\textwidth]{Problem/AIME_1983_I_15_Problem_Figure_1.png}
                    \end{figure}
                    


 
 \textbf{Answer: }175



\textbf{Solution 1}

As with some of the other solutions, we analyze this with a locus—but a different one.  We'll consider: given a point $P$ and a line $\ell,$ what is the set of points $X$ such that the midpoint of $PX$ lies on line $\ell$?  The answer to this question is: a line $m$ parallel to $\ell$, such that $m$ and $P$ are (1) on opposite sides of $\ell$. and (2) at the same distance from $\ell$.

 
                    \begin{figure}[H]
                        \centering
                        \includegraphics[width=0.5\textwidth]{Solution/AIME_1983_I_15_Solution_Figure_1.png}
                    \end{figure}
                    Applied to this problem, this means that $D$ is the only point that lies on both (1) the given circle, and (2) the line through $D$ parallel to $BC$.  This means that $BC$ is parallel to the tangent to the given circle at $D$.

 If we take $O$ to be the center of the given circle, then this means that $OD$ is perpendicular to $BC$.  Let $M$ be the midpoint of chord $BC,$ and let $N$ be the intersection of $OD$ with the line through $A$ parallel to $BC$.

 
                    \begin{figure}[H]
                        \centering
                        \includegraphics[width=0.5\textwidth]{Solution/AIME_1983_I_15_Solution_Figure_1.png}
                    \end{figure}
                    Since $BC = 6,$ we know that $BM = 3$; since $OB$ (a radius of the circle) is 5, we can conclude that $\triangle BMO$ is a 3-4-5 right triangle.  Since $D$ and line $AN$ are equidistant from line $BC,$ we know that $MN = 1$, and thus $ON = 3$.  This makes $\triangle ANO$ also a 3-4-5 right triangle.

 We're looking for $\sin \angle AOB$, and we can find that using the \href{https://artofproblemsolving.com/wiki/index.php?title=Trigonometric\_identities\#Angle\_addition\_identities}{ angle subtraction formula} for sine.  We have
 \begin{align*} \sin \angle AOB &= \sin(\angle AOM - \angle BOM) \\ &= \sin \angle AOM \cos \angle BOM - \cos \angle AOM \sin \angle BOM \\ &= \frac{4}{5} \cdot \frac{4}{5}  - \frac{3}{5} \cdot \frac{3}{5} \\ &= \frac{16 - 9}{25} = \frac{7}{25}. \end{align*}
This is in lowest terms, so our answer is $mn = 7 \cdot 25 = 175$.

 





\textbf{Solution 2}

-Credit to Adamz for diagram-
                    \begin{figure}[H]
                        \centering
                        \includegraphics[width=0.5\textwidth]{Solution/AIME_1983_I_15_Solution_Figure_1.png}
                    \end{figure}
                    Let $A$ be any fixed point on circle $O$, and let $AD$ be a chord of circle $O$. The \href{https://artofproblemsolving.com/wiki/index.php?title=Locus}{locus} of midpoints $N$ of the chord $AD$ is a circle $P$, with diameter $AO$. Generally, the circle $P$ can intersect the chord $BC$ at two points, one point, or they may not have a point of intersection. By the problem condition, however, the circle $P$ is tangent to $BC$ at point $N$. 

 Let $M$ be the midpoint of the chord $BC$. From right triangle $OMB$, we have $OM = \sqrt{OB^2 - BM^2} =4$. This gives $\tan \angle BOM = \frac{BM}{OM} = \frac 3 4$.

 Notice that the distance $OM$ equals $PN + PO \cos \angle AOM = r(1 + \cos \angle AOM)$, where $r$ is the radius of circle $P$. 

 Hence 
 \[\cos \angle AOM = \frac{OM}{r} - 1 = \frac{2OM}{R} - 1 = \frac 8 5 - 1 = \frac 3 5\] (where $R$ represents the radius, $5$, of the large circle given in the question). Therefore, since $\angle AOM$ is clearly acute, we see that 
 \[\tan \angle AOM =\frac{\sqrt{1 - \cos^2 \angle AOM}}{\cos \angle AOM} = \frac{\sqrt{5^2 - 3^2}}{3} = \frac 4 3\]

 Next, notice that $\angle AOB = \angle AOM - \angle BOM$. We can therefore apply the subtraction formula for $\tan$ to obtain 
 \[\tan \angle AOB =\frac{\tan \angle AOM - \tan \angle BOM}{1 + \tan \angle AOM \cdot \tan \angle BOM} =\frac{\frac 4 3 - \frac 3 4}{1 + \frac 4 3 \cdot \frac 3 4} = \frac{7}{24}\] It follows that $\sin \angle AOB =\frac{7}{\sqrt{7^2+24^2}} = \frac{7}{25}$, such that the answer is $7 \cdot 25=\boxed{175}$.

 





\textbf{Solution 3}

This solution, while similar to Solution 2, is arguably more motivated and less contrived.

 Firstly, we note the statement in the problem that "$AD$ is the only chord starting at $A$ and bisected by $BC$" – what is its significance? What is the criterion for this statement to be true?

 We consider the locus of midpoints of the chords from $A$. It is well-known that this is the circle with diameter $AO$, where $O$ is the center of the circle. The proof is simple: every midpoint of a chord is a dilation of the endpoint with scale factor $\frac{1}{2}$ and center $A$. Thus, the locus is the result of the dilation with scale factor $\frac{1}{2}$ and centre $A$ of circle $O$. Let the center of this circle be $P$.

 Now, $AD$ is bisected by $BC$ if they cross at some point $N$ on the circle. Moreover, since $AD$ is the only chord, $BC$ must be tangent to the circle $P$.

 The rest of this problem is straightforward.

 Our goal is to find $\sin \angle AOB = \sin{\left(\angle AOM - \angle BOM\right)}$, where $M$ is the midpoint of $BC$. We have $BM=3$ and $OM=4$.Let $R$ be the projection of $A$ onto $OM$, and similarly let $Q$ be the projection of $P$ onto $OM$. Then it remains to find $AR$ so that we can use the addition formula for $\sin$.

 As $PN$ is a radius of circle $P$, $PN=2.5$, and similarly, $PO=2.5$. Since $OM=4$, we have $OQ=OM-QM=OM-PN=4-2.5=1.5$. Thus $PQ=\sqrt{2.5^2-1.5^2}=2$.

 Further, we see that $\triangle OAR$ is a dilation of $\triangle OPQ$ about center $O$ with scale factor $2$, so $AR=2PQ=4$.

 Lastly, we apply the formula:
 \[\sin{\left(\angle AOM - \angle BOM\right)} = \sin \angle AOM \cos \angle BOM - \sin \angle BOM \cos \angle AOM = \left(\frac{4}{5}\right)\left(\frac{4}{5}\right)-\left(\frac{3}{5}\right)\left(\frac{3}{5}\right)=\frac{7}{25}\]Thus the answer is $7\cdot25=\boxed{175}$.

 





\textbf{Solution 4 (coordinate geometry)}


                    \begin{figure}[H]
                        \centering
                        \includegraphics[width=0.5\textwidth]{Solution/AIME_1983_I_15_Solution_Figure_1.png}
                    \end{figure}
                    

 Let the circle have equation $x^2 + y^2 = 25$, with centre $O(0,0)$. Since $BC=6$, we can calculate (by the Pythagorean Theorem) that the distance from $O$ to the line $BC$ is $4$. Therefore, we can let $B=(3,4)$ and $C=(-3,4)$. Now, assume that $A$ is any point on the major arc BC, and $D$ any point on the minor arc BC. We can write $A=(5 \cos \alpha, 5 \sin \alpha)$, where $\alpha$ is the angle measured from the positive $x$ axis to the ray $OA$. It will also be convenient to define $\angle XOB = \alpha_0$. 

 Firstly, since $B$ must lie in the minor arc $AD$, we see that $\alpha \in \left(-\frac{\pi}{2}-\alpha_0,\alpha_0\right)$. However, since the midpoint of $AD$ must lie on $BC$, and the highest possible $y$-coordinate of $D$ is $5$, we see that the $y$-coordinate cannot be lower than $3$, that is, $\alpha \in \left[\sin^{-1}\frac{3}{5},\alpha_0\right)$.

 Secondly, there is a theorem that says that, in a circle, if a chord is bisected by a radius, then they must be perpendicular. Therefore, suppose that $P$ is the intersection point of $AD$ and $BC$, so that by the theorem, $OP$ is perpendicular to $AD$. So, if $AD$ is the only chord starting at $A$ which is bisected by $BC$, this means that $P$ is the only point on the chord $BC$ such that $OP$ is perpendicular to $AD$. Now suppose that $P=(p,4)$, where $p \in (-3,3)$. The fact that $OP$ must be perpendicular to $AD$ is equivalent to the following equation:

 
 \[-1 = \left(\text{slope of } OP\right)\left(\text{slope of } AP\right)\]which becomes
 \[-1 = \frac{4}{p} \cdot \frac{5\sin \alpha - 4}{5\cos \alpha - p}\]

 This rearranges to

 
 \[p^2 - (5\cos \alpha)p + 16 - 20 \sin \alpha = 0\]

 Given that this equation must have only one real root $p\in (-3,3)$, we study the following function:

 
 \[f(x) =  x^2 - (5\cos \alpha)x + 16 - 20 \sin \alpha\]

 First, by the fact that the equation $f(x)=0$ has real solutions, its discriminant $\Delta$ must be non-negative, so we calculate

 
 \[\begin{split}\Delta & = (5\cos \alpha)^2 - 4(16-20\sin \alpha) \\ & =  25 (1- \sin^2 \alpha)  - 64 + 80 \sin \alpha \\ & = -25 \sin^2 \alpha + 80\sin \alpha - 39 \\ & =  (13 - 5\sin \alpha)(5\sin \alpha - 3)\end{split}\]

 It is obvious that this is in fact non-negative. If it is actually zero, then $\sin \alpha = \frac{3}{5}$, and $\cos \alpha = \frac{4}{5}$. In this case, $p = \frac{5\cos \alpha}{2} = 2 \in (-3,3)$, so we have found a possible solution. We thus calculate $\sin(\text{central angle of minor arc } AB) = \sin (\alpha_0 - \alpha) = \frac{4}{5}\cdot \frac{4}{5} - \frac{3}{5} \cdot \frac{3}{5} = \frac{7}{25}$ by the subtraction formula for $\sin$. This means that the answer is $7 \cdot 25 = 175$.

 \textbf{Addendum to Solution 4} Since this is an AIME problem, we can assume that we are done since we have found one possible case. However, in a full-solution contest, we could not assume that the answer is unique, but would need to prove that this is the unique solution. This can be proven as follows.

 Suppose that $\Delta > 0$, which would mean that there could be two real roots of $f(x)$, one lying in the interval $(-3,3)$, and another outside of it. We also see, by \href{https://artofproblemsolving.com/wiki/index.php?title=Vieta\%27s\_Formulas}{Vieta's Formulas}, that the average of the two roots is $\frac{5\cos \alpha}{2}$, which is non-negative, so the root outside of $(-3,3)$ must be no less than $3$. By considering the graph of $y=f(x)$, which is a "U-shaped" parabola, it is now evident that $f(-3) > 0$ and $f(3)\leq 0$. We can just use the second inequality:

 
 \[0 \geq  f(3) =  25  - 15\cos \alpha - 20 \sin \alpha\]so
 \[3\cos \alpha + 4 \sin \alpha  \geq 5\]

 The only way for this inequality to be satisfied is when $A=B$ (by applying the Cauchy-Schwarz inequality, or just plotting the line $3x+4y=5$ to see that point $A$ can't be above this line), which does not make sense in the original problem statement. (For it would mean that the point $A$ lies in the half-plane above the line $3x+4y=5$, inclusive, and the half-plane below the line $-3x+4y=5$, exclusive. This can be seen to be impossible by drawing the lines and observing that the intersection of the two half-planes does not share any point with the circle.)

 





\textbf{Solution 5}

Let the center of the circle be $O$. Fix $B,C,$ and $A$. Then, as $D$ moves around the circle, the locus of the midpoints of $AD$ is clearly a circle. Since the problems gives that $AD$ is the only chord starting at $A$ bisected by $BC$, it follows that the circle with diameter $DO$ and $AO$ is tangent to $BC$.

 Now, let the intersection of $BC$ and $AD$ be $E$ and let the midpoint of $AO$ (the center of the circle tangent to $BC$ that we described beforehand) be $F$. Drop the altitude from $O$ to $BC$ and call its intersection with $BC$ $K$. Drop the perpendicular from $F$ to $KO$ and call its intersection with $KO$ $L$. Clearly, $KO = \sqrt{OC^2-KC^2} = \sqrt{5^2-3^2} = 4$ and since $EF$ is radius, it equals $\frac{5}{2}$. The same applies for $FO$, which also equals $\frac{5}{2}$. By the Pythagorean theorem, we deduce that $FL = 2$, so $EK = 2$. This is very important information! Now we know that $BE = 1$, so by Power of a Point, $AE = ED = \sqrt{5}$.

 We’re almost there! Since by the Pythagorean theorem, $ED^2 + EO^2 = 25$, we deduce that $EO = 2\sqrt{5}$. $EC=OC=5$, so $\sin (CEO) = \frac{2\sqrt{5}}{5}$. Furthermore, since $\sin (CEO) = \cos(DEC)$, we know that $\cos (DEC) = \frac{2\sqrt{5}}{5}$. By the law of cosines,
 \[DC^2 = (\sqrt{5})^2 + 5^2 -10\sqrt{5} \cdot \frac{2\sqrt{5}}{5} = 10\]Therefore, $DC = \sqrt{10} \Longleftrightarrow BA = \sqrt{2}$. Now, drop the altitude from $O$ to $BA$ and call its intersection with $BA$ $Z$. Then, by the Pythagorean theorem, $OZ = \frac{7\sqrt{2}}{2}$. Thus, $\sin (BOZ) = \frac{\sqrt{2}}{10}$ and $\cos (BOZ) = \frac{7\sqrt{2}}{10}$. As a result, $\sin (BOA) = \sin (2 BOZ) = 2\sin(BOZ)\cos(BOZ) = \frac{7}{25}$. $7 \cdot 25 = \boxed{175}$.

 





\textbf{Solution 6}


                    \begin{figure}[H]
                        \centering
                        \includegraphics[width=0.5\textwidth]{Solution/AIME_1983_I_15_Solution_Figure_1.png}
                    \end{figure}
                    Let I be the intersection of AD and BC.

 Lemma: $AI = ID$ if and only if $\angle AIO = 90$.

 Proof: If AI = ID, we get AO = OD, and thus IO is a perpendicular bisector of AD. If $\angle AIO = 90$, We can get $\triangle AIO \cong \triangle OID$

 Let  be this the circle with diameter AO.

 Thus, we get $\angle AIO = 90$, implying I must lie on $\omega$. I also must lie on BC. If both of these conditions are satisfied, The extension of AI intersecting with circle O will create a point D such that AI = ID. Thus, since AD is the only chord starting at A, there must be only 1 possible I, implying the circle with diameter AO is tangent to BC.

 Now, create a coordinate system such that the origin is O and the y-axis is perpendicular to BC. The positive x direction is from B to C.

 Let Z be (0,5).Let Y be (-5,0).Let X be the center of $\omega$. Since $\omega$'s radius is $\frac{5}{2}$, the altitude from X to OY is 1.5, since the altitude from I to OY is 4. XO is 2.5, so $sin(XOY) = sin(AOY) = \frac{3}{5}$. $sin(BOZ) =  \frac{3}{5}$. If we let $sin(\theta) = \frac{3}{5}$, we can find that what we are looking for is $sin(90 - 2\theta)$, which we can evaluate and get $\frac{7}{25} \implies \boxed{175}$

 -Alexlikemath

 





\textbf{Solution 7 (No Trig)}

Let $O$ be the center of the circle. The locus of midpoints of chords with $A$ as a endpoint is a circle with diameter $\overline{AO}$. Additionally, this circle must be tangent to $\overline{BC}$. Let the center of this circle be $P$. Let $M$ be the midpoint of $BC$, $N$ be the foot of the perpendicular from $P$ to $\overline{BM}$, and $K$ be the foot of the perpendicular from $B$ to $\overline{AP}$. Let $x=BK$.

 From right triangle $BKO$, we get $KO = \sqrt{25-x^2}$. Thus, $KP = \sqrt{25-x^2}-\frac52$.

 Since $BO = 5$, $BM = 3$, and $\angle BMO$ is right, $MO=4$. From quadrilateral $MNPO$, we get $MN = \sqrt{PO^2 - (MO - NP)^2} = \sqrt{(5/2)^2 - (4 - 5/2)^2} = \sqrt{(5/2)^2 - (3/2)^2} = 2$. Thus, $BN = 1$.

 Since angles $BNP$ and $BKP$ are right, we get 
 \[BK^2+KP^2 = BN^2 + NP^2 \implies x^2 + \left(\sqrt{25-x^2}-\frac52\right)^2 = \left(\frac52\right)^2 + 1\]
 \[25 - 5\sqrt{25-x^2} = 1\]
 \[5\sqrt{25-x^2} = 24\]
 \[25(25-x^2) = 24^2\]
 \[25x^2 = 25^2 - 24^2 = 49\]
 \[x = \frac75\]Thus, $\sin \angle AOB = \frac{x}{5} = \frac{7}{25}\implies \boxed{175}$.

 





\textbf{Solution 8 (Similar to 1, but with coordinates)}

Let the center of the circle be O. O is at (0,0).Rotate the circle so that the line BC has slope 0, and so that C is in the 1st quadrant. Since BC = 6, and B can be reflected across the y axis to become C, we can say that B is on x=-3, and C is on x=3. Since it must lie on the circle $x^2+y^2=25$,B = (-3,4) and C = (3,4). So BC is on the line y=4.

 Let A be at $(x,\sqrt{25-x^2})$. Let D be at $(z,\sqrt{25-z^2})$.Notice that the midpoint of AD lies on BC. Since BC is on the line y=4, using the midpoint formula we can say, $\sqrt{25-x^2} + \sqrt{25-z^2} = 8$.We treat x like a constant, since it is determined by where A lies on the circle.Notice that if the above equation is true for z=p, it is also true for z=-p. But this is impossible because the problem states that AD is the only chord starting at A which is bisected by BC. This problem is solved if z=-z=0. Thus, z=0, and D=(0,5). plugging in z=0 into $\sqrt{25-x^2} + \sqrt{25-z^2} = 8$, we find x=+-4. Since AB is a minor arc, we assume x=-4. A is therefore on (-4,3) and B on (-3,4).

 We can use complex numbers to find the angle between A and B, because arguments subtract when you divide complex numbers. $\frac{-4+3i}{-3+4i} = \frac{24+7i}{25}$ Which gives us $\frac{7}{25}$. $7 \cdot 25 = \boxed{175}$

 




\end{enumerate}